% AgentMPI: A Message Passing Interface for Multi-Agent Harness Development
\documentclass[10pt,twocolumn]{article}

\usepackage[margin=0.75in]{geometry}
\usepackage[T1]{fontenc}
\usepackage[utf8]{inputenc}
\usepackage{lmodern}
\usepackage{microtype}
\usepackage{amsmath,amssymb}
\usepackage{booktabs}
\usepackage{multirow}
\usepackage{graphicx}
\usepackage{xcolor}
\usepackage{listings}
\usepackage[hidelinks]{hyperref}
\usepackage[capitalise,noabbrev]{cleveref}
\usepackage{balance}
\usepackage{flushend}

\lstset{
  basicstyle=\ttfamily\scriptsize,
  keywordstyle=\bfseries,
  commentstyle=\itshape\color{gray!70!black},
  columns=fullflexible,
  breaklines=true,
  frame=none,
  aboveskip=4pt, belowskip=4pt,
  showstringspaces=false,
  language=Python,
  morekeywords={def,for,if,else,with,return,while,in,not,and,or,None,True,False}
}

\usepackage{amsthm}
\theoremstyle{definition}
\newtheorem{definition}{Definition}
\theoremstyle{plain}
\newtheorem{proposition}{Proposition}

\newcommand{\ampi}{\textsc{AgentMPI}}
\newcommand{\code}[1]{\texttt{\small #1}}
\newcommand{\NAgentCallsTranslation}{112}
\newcommand{\NAgentCallsSoftware}{16}
\newcommand{\NAgentCallsMicro}{0}
\newcommand{\NAgentCallsTotal}{128}
\newcommand{\NTranslationConfigs}{6}
\newcommand{\NEagerLimit}{4,096}
\newcommand{\NModelAgreement}{100}
\newcommand{\NModelConfigs}{278}
\newcommand{\NModelDepthAgreement}{100}
\newcommand{\NUslFit}{$\sigma=0.220$, $\kappa=0.0000$ ($R^2=0.873$)}


\title{\textbf{\ampi{}: A Message Passing Interface for\\Multi-Agent Harness Development}}
\author{}
\date{}

\begin{document}
\maketitle

\begin{abstract}
Multi-agent systems built on large language models are written the way parallel
programs were written before 1994: every project invents its own communication
mechanism, and none of them compose. We argue that the field needs what MPI gave
parallel computing --- not a framework, but a \emph{library-level protocol} with
scoped communication contexts, typed transfers, collective operations with
explicit algorithms, one-sided shared state with epochs, and failure mitigation
--- and we show that MPI's abstractions transfer with three specific, load-bearing
modifications.

First, the receiving agent's context window is a scarce, overcommitted buffer, so
the eager/rendezvous distinction that MPI hides as an implementation detail must
become semantically visible; making it visible turns MPI's notion of an
\emph{unsafe program} into a precise and checkable notion of \emph{context
safety}. Second, reduction operators implemented by a language model are not
associative and are \emph{lossy in proportion to the depth of the reduction},
which makes collective-algorithm selection a quality decision rather than only a
cost decision --- with the surprising consequence that the logarithmic-depth tree
wins on fidelity as well as latency, provided broadcasts forward immutable
content-addressed handles rather than regenerated text. Third, the dominant agent
failure is not the crash MPI's fault-tolerance work assumes but a well-formed,
confident, wrong answer; this is the analogue of silent data corruption, for which
the correct response is verification rather than checkpointing.

We specify \ampi{}, implement it in 6\,900 lines of Python with all collectives
built over point-to-point so that their structure is observable, and validate the
implementation against closed-form cost models on \NModelConfigs{}
collective configurations with \NModelAgreement\%
agreement on message counts and fold depths. We evaluate with real agent
populations on two tasks chosen to bracket the coupling spectrum: parallel
translation of a book, where the dependence is weak but global, and collaborative
implementation of a specified software system against a fixed acceptance suite,
where it is not.
\end{abstract}

\section{Introduction}
\label{sec:intro}

In 1993 there was roughly one message-passing library per parallel machine. PVM,
p4, Chameleon, PARMACS, Express, Zipcode, NX/2 and CMMD each solved the same
problem differently, and a program written for one had to be rewritten for the
next~\cite{walker92,geist94-pvm,butler94-p4,calkin94-parmacs,express92,skjellum90}.
The MPI Forum's response was not a better library but a \emph{specification}: a
portable interface that vendors could implement well and that application and
library authors could target once~\cite{mpiforum94,dongarra96}. Two decades of
scientific software rest on that decision.

Multi-agent systems built on language models are at the same point.
AutoGen, MetaGPT, ChatDev, CAMEL, LangGraph, CrewAI and their successors each
supply a coordination model, and a harness written against one cannot be moved to
another or composed with
another~\cite{wu24autogen,hong24metagpt,qian24chatdev,li23camel,langgraph,crewai}.
The protocols that have emerged --- MCP, A2A, ACP, ANP --- are, on inspection,
about a different problem: MCP standardises agent-to-\emph{tool} access, and A2A
standardises task delegation \emph{across organisational
boundaries}~\cite{mcpspec2026,a2adefs,anpwp,coral}. Neither provides what an
author of a multi-agent \emph{program} needs: a way to name a subset of the
population and communicate with it privately, a collective operation, a shared
region with defined concurrency semantics, or a defined behaviour when a
participant disappears.

The consequences are documented, and --- more usefully --- they are documented in
terms a message-passing specification already defines. The MAST taxonomy of failed
multi-agent traces attributes failures to specification and inter-agent
coordination rather than to model capability~\cite{mast,mastneurips}, and roughly
42\% of its observed failures fall in categories that name MPI mechanisms directly:
step repetition (17.1\%, an idempotence property), termination-unaware and
premature termination (17.6\%, a barrier and agreement property), history loss and
resets (5.7\%, a durable log), and withheld or ignored input (1.8\%, delivery
semantics). We are careful about the complement: reasoning--action mismatch (14.0\%)
and disobeying the task specification (11.0\%) are \emph{not} protocol-addressable,
and no communication layer will fix them.

Independent evidence points the same way. Mieczkowski et al.\ analyse language-model
teams explicitly as distributed systems and find Amdahl's law bounding team
speedup, decentralised teams reaching a median speedup below unity --- slower than a
single agent --- and concurrent-write conflicts producing several times more failed
tests than a coordinated arrangement~\cite{lmteams26}. Practitioners report the same
outcome and attribute it to context fragmentation and information lost between
agents~\cite{cognitiondont,anthropicmas}. Long-context degradation and multi-turn
drift are measured and substantial~\cite{liu24lostinmiddle,labanlost,chromarot},
and dense communication topologies are shown to \emph{cause} premature convergence,
which argues for scoped sharing as a per-operation choice rather than for more
sharing~\cite{diversitycollapse}.

These are not model problems. They are the problems a communication protocol exists
to solve, and in 1993 the parallel-computing community had all of them too.

\subsection{This paper}

We ask what MPI actually contributes, transplant what transfers, and are explicit
about what does not.

What transfers, and matters more than expected:
\begin{itemize}\itemsep2pt
\item \textbf{The communicator}, MPI's pair of a process group and an opaque
communication context. It exists because a library and its caller could otherwise
collide on an integer tag. The agent version of that collision is not
hypothetical: it is a reviewer's critique reaching the wrong author, and the
single-shared-conversation architecture makes it certain (\cref{sec:model}).
\item \textbf{Derived datatypes.} Their under-appreciated purpose is describing a
\emph{projection} of data so it can be communicated without being copied. In the
agent setting this is the mechanism by which a rank receives what it needs out of
a 400k-token artifact rather than the artifact (\cref{sec:ptp}).
\item \textbf{Neighbourhood collectives on a virtual topology.} An all-to-all
conversation over $p$ agents costs $\Theta(p^2)$ messages; over a ring or review
graph the same information costs $\Theta(p)$ (\cref{sec:coll}).
\item \textbf{One-sided operations with epochs and the \code{SEPARATE} memory
model.} Every multi-agent system grows a shared scratchpad and gets its
concurrency wrong the same way. MPI-3 supplies both the vocabulary and the fix
(\cref{sec:rma}).
\item \textbf{ULFM's \code{revoke}/\code{shrink}/\code{agree}}, whose least
obvious member, \code{revoke}, is the one that makes recovery possible at all
(\cref{sec:ft}).
\end{itemize}

What must change, and why the changes are the technical contribution:

\paragraph{C1: context is the scarce buffer, and transport mode must be visible.}
MPI chooses between eager and rendezvous transfer invisibly, because both deliver
the same bytes. For an agent rank they do not have the same effect: an eager
payload enters the agent's context window and a rendezvous envelope does not.
Making the mode explicit lets us restate MPI's discipline of the \emph{safe
program} --- one that does not depend on buffering --- as \emph{context safety},
and to show experimentally that the same ring-exchange program fails above a
payload threshold under eager transport and completes at every size under
rendezvous (\cref{sec:ptp}, \cref{tab:transport}).

\paragraph{C2: reduction is lossy in depth, so algorithm choice is a quality
decision --- and the optimal tree is not binary.} MPI requires reduction operators
to be associative and then reorders freely; the resulting non-reproducibility of
floating-point \code{MPI\_SUM} is tolerated because it is bounded by rounding
error~\cite{Demmel15,IntelCNR}. An operator implemented by a language model is not
associative and is lossy, and its loss compounds with the \emph{depth} of the
reduction rather than the number of applications. We make algebraic strength part
of an operator's declaration, use it to constrain algorithm selection, and report
fold depth from every reduction.

Two consequences follow, and the second was not what we expected. First, a binomial
tree at depth $\lceil\log_2 p\rceil$ beats a serial chain at depth $p-1$ on
fidelity \emph{and} latency --- but only because broadcasts forward
content-addressed handles, so relaying cannot corrupt. Second, and more
interestingly: \textbf{binary is the wrong arity.} Every logarithmic collective in
MPI descends from processes being \emph{single-ported} --- a process participates in
one transfer at a time, so a wider tree node buys nothing. An agent rank has no
port count. A prompt carries eight reports as easily as two, and merging eight in
one call is one lossy re-encoding rather than seven. The binding constraint becomes
the receiving rank's context budget, giving an optimal fan-in
$k^{*}=\lfloor C(1-h)/s\rfloor$ and a fold depth of $\lceil\log_{k^{*}} p\rceil$.
At $p=64$, widening from $k{=}2$ to $k{=}8$ takes the fold depth from 6 to 2
\emph{at an identical message count}, since every non-root rank sends exactly once
whatever the arity. So the agent setting should prefer the \emph{widest} feasible
tree where MPI prefers the narrowest (\cref{sec:coll}, \cref{tab:fanin}).

\paragraph{C3: the dominant failure is silent corruption, not crash.} A rank that
returns a confident wrong answer is invisible to timeouts, retries and schema
checks. It is the analogue of silent data corruption, and HPC's answer to that is
not checkpoint/restart --- which preserves the corruption faithfully --- but
algorithm-based fault tolerance~\cite{huang84abft}. We give a six-class failure
model, pair each class with the only detector that can see it, and make
verification a first-class protocol concept (\cref{sec:ft}).

\paragraph{C4: protocol in the harness, not in the prompt.} Because an agent's
working memory is lossy and its instances are ephemeral, no protocol state may
live in the agent. An agent holds only handles; every \ampi{} call is made by
trusted host-side code, and the model is invoked as a kernel that transforms
artifacts (\cref{sec:model}).

\subsection{Contributions}
\begin{enumerate}\itemsep2pt
\item A specification of \ampi{} (\cref{sec:model}--\cref{sec:cost}), written in
the MPI standard's style with explicit rationale, covering communicators,
point-to-point transfer with contracts and views, twelve collectives with
declared algorithms, one-sided windows with epochs and leases, a failure model
with mitigations, and a three-axis cost model.
\item A reference implementation in which every collective is built over
point-to-point, so message counts, rounds and fold depths are observable, together
with closed-form cost formulas for all 27 implemented (collective, algorithm)
pairs and a validation showing exact agreement on
\NModelConfigs{} configurations.
\item Two end-to-end experiments with real agent populations that bracket the
coupling spectrum, plus microbenchmarks establishing the calibration constants,
the context-safety threshold, fault-recovery costs and reduction fidelity.
\item A calibrated discrete-event simulator, validated against measured runs, that
extrapolates collective cost to populations larger than can be run.
\end{enumerate}

\section{Background: what MPI contributes}
\label{sec:background}

\subsection{The standardisation decision}

MPI-1 was drafted between late 1992 and May 1994 by a forum of vendors, national
laboratories and universities~\cite{mpiforum93,mpiforum94}. It defined an
\emph{interface}, not an implementation, and it deliberately omitted much: no
explicit shared memory, no parallel I/O, no dynamic process creation, and no fault
tolerance. Those arrived in MPI-2 (one-sided operations,
\code{MPI\_Comm\_spawn}, MPI-IO)~\cite{mpi20}, MPI-3 (nonblocking and
neighbourhood collectives, a redesigned one-sided interface, shared-memory
windows, matched probes)~\cite{mpi31,Hoefler-MPI3-overview,hoefler15-rma,MPI31-mprobe},
and MPI-4 (persistent collectives, partitioned communication,
sessions)~\cite{mpi40}. The omissions are as instructive as the inclusions: by
declining to be a language, to guarantee buffering, or to handle faults, MPI-1
stayed small enough that every vendor implemented it within two
years~\cite{snir96-whybig,snir96-buffering}.

\subsection{The four ideas worth stealing}

\paragraph{Communicators separate the message space from the address space.}
A communicator is a group plus an opaque context, and \code{MPI\_Comm\_dup} exists
to manufacture a fresh context so that a library can communicate without
interference from its caller~\cite{mpiforum94,gropp99}. This is what makes
parallel libraries composable, and composability is exactly what current agent
frameworks lack.

\paragraph{Datatypes do two jobs.} The first is matching: a type signature on both
sides, checked at runtime, converting silent corruption into a loud error. The
second, and the one that motivated the design, is describing non-contiguous data
so that a boundary can be communicated without packing a
copy~\cite{MPI-datatypes,Traff99,Reussner00}.

\paragraph{Collectives are algorithms, and the algorithm matters.} MPI's
collective layer is built over point-to-point, and decades of work
characterise which algorithm wins in which regime under the Hockney
$\alpha$--$\beta$ model and its refinements
LogP/LogGP~\cite{Hockney94,Culler93,Alexandrov97,ThakurRabenseifnerGropp05,Rabenseifner04,Bruck97,Chan07,Sanders09}.
The results are quantitative: binomial-tree broadcast at
$\lceil\log_2 p\rceil(\alpha+n\beta)$ versus flat broadcast at
$(p-1)(\alpha+n\beta)$; Bruck's algorithm trading volume for rounds; Rabenseifner's
reduce-scatter-plus-allgather allreduce.

\paragraph{Safety is a property of programs, not implementations.} MPI declines to
guarantee that a program depending on buffering will work, because such programs
run on one implementation and deadlock on another~\cite{snir96-buffering}. Naming
the hazard and pushing it to the programmer, with \code{MPI\_Sendrecv} as the
remedy, is a design pattern we reuse directly.

\subsection{What MPI gets wrong for this setting}

MPI's fault-tolerance posture is that after an error its state is undefined and
the default handler aborts. FT-MPI~\cite{fagg00-ftmpi} and then
ULFM~\cite{bouteiller22} supply mitigation primitives, but they assume fail-stop
processes. MPI also has no notion of a receive buffer that cannot be resized, no
notion of an operator that is only approximately associative, and no vocabulary
for a participant that answers confidently and wrongly. Those gaps are what
\cref{sec:ptp}--\cref{sec:ft} fill.

\section{The \ampi{} model}
\label{sec:model}

\subsection{Ranks are roles; incarnations are sessions}

A job has $p$ ranks, fixed at creation. A rank is a \emph{durable role} defined by
its number, mailbox, context budget and accounting. Its embodiment is an
\emph{incarnation}: one agent session bound to it for a bounded interval under a
lease. The mapping is one-to-many over time and exactly-one at any instant.

In MPI a rank \emph{is} a process, and the identification is harmless because MPI
processes live for the whole job. Agent sessions do not: they end at a timeout, at
context exhaustion, or at the host's discretion. Tying rank identity to session
identity would force a communicator rebuild on every such event and invalidate the
harness's mapping from rank to work --- and for agents that mapping is expensive,
because ``rank 7 owns the parser'' is baked into prompts and artifacts. Separating
them makes session termination ordinary and recoverable: the rank persists, its
durable mailbox persists, and a fresh incarnation resumes it. This is the virtual
actor discipline~\cite{bernstein14orleans} grafted onto MPI's static rank space.

\subsection{Artifacts are immutable and content-addressed}

Everything crossing a communicator is an immutable byte string addressed by the
SHA-256 of its canonical serialisation. Three consequences are normative, and the
first is the one the design depends on:

\begin{quote}
\emph{A tree broadcast forwards digests, never regenerated content.}
\end{quote}

The obvious agent implementation of tree broadcast --- ``tell your children what I
told you'' --- is a game of telephone whose corruption grows with depth, which
would confine the protocol to flat, root-centred patterns and make the root a
bottleneck at every scale. Forwarding handles makes a
$\lceil\log_2 p\rceil$-depth broadcast deliver byte-identical content to every
rank. It also draws the line that \cref{sec:coll} rests on: \textbf{broadcast can
be lossless because it does not transform; reduction cannot be, because it must.}

The second consequence is that admitting the same digest twice to a rank's context
costs the budget once, which matters because broadcast-then-echo is pervasive. The
third is that the event log plus the artifact store suffice to replay the protocol
without invoking any model.

\subsection{Agents hold only handles}

An agent is never required to remember protocol state. It receives handles: a rank
number, a communicator id, a window slot, a digest, an invocation id. We take this
further than seems necessary --- where an agent must issue a command, the runtime
emits the \emph{exact command string} rather than the parameters from which it could
be constructed. Protocol violation in agent populations is dominated by recall
failures, so a surface requiring recognition rather than recall is materially more
reliable.

We can report a sharp, accidental experiment on precisely this point, because our
own implementation got it wrong. In the first campaign the emitted submission
command was itself malformed --- an option placed after the subcommand it belonged
before --- and all nine workers in the population hit the parse error. All nine
diagnosed it, reordered the flag, and completed their submissions, several
independently reporting the defect back as something the runtime should fix. The
population was robust; the mechanism was not. Two conclusions, and the second is the
uncomfortable one. A generated command handed to an agent is executable code and
must be tested as such, so the emitted strings are now round-tripped through the
real argument parser in the test suite. And the recovery is not reassurance: these
workers had shell access, could read an error, and were willing to retry. A worker
that treated submission as fire-and-forget would have silently discarded work that
was already correct --- which is exactly the class of failure the handle discipline
exists to prevent, arriving through the mechanism meant to prevent it.

\subsection{Protocol in the harness, not in the prompt}

\ampi{} supports two harness styles. In the \emph{harness-side} (SPMD) style the
author writes a \code{rank\_main} executed once per rank by trusted host-side code;
every \ampi{} call is made by that code and the model is invoked only through
\code{comm.agent}. In the \emph{agent-side} style the model itself issues protocol
operations through a command-line interface.

\begin{lstlisting}
def rank_main(comm):
    # decompose: variable-sized blocks, by handle
    mine = comm.scatterv(work if comm.rank == 0 else None, root=0)
    # agree a glossary: exact, idempotent, tree-parallel
    terms = comm.agent(extract_prompt(mine), contract=TERMSHEET)
    glossary = comm.allreduce(terms, ampi.UNION)
    out = comm.agent(translate_prompt(mine, glossary),
                     contract=TRANSLATION)
    # reconcile seams with immediate neighbours only
    topo = ampi.cart_create(comm, dims=[comm.size], periods=[False])
    left, right = ampi.halo_exchange(topo, head(out), tail(out))
    out = comm.agent(revise_prompt(out, left, right))
    # collect by handle: the root never admits p*n tokens
    return comm.gather(out, root=0, mode=Mode.RENDEZVOUS)
\end{lstlisting}

The distinction is substantive. In the agent-side style, protocol conformance is a
property of model behaviour, and a rank that forgets to enter a barrier does not
merely give a worse answer --- it prevents the population from making progress.
Confining the protocol to host-side code makes conformance a property of the
runtime. We support the agent-side style because it is sometimes the only option,
and because the difference is measurable.

\subsection{Protocol state lives outside the agents}

MPI keeps protocol state in process memory, which is exactly as reliable as the
process. For agent ranks that assumption fails twice over: the agent's working
memory is a lossy context window, and the agent instance routinely disappears. So
\ampi{} keeps every byte of protocol state --- posted receives, unexpected
messages, window contents, communicator contexts, leases --- in an external store
with ACID semantics, and a rank can be killed between any two calls and rebuilt
from that store alone. The reference fabric is SQLite in WAL mode plus a
content-addressed blob store; that is a correctness choice, not a performance one,
and \cref{sec:eval-micro} shows the per-operation protocol cost is roughly four
orders of magnitude below an agent invocation, so it is free.

\section{Point-to-point and context safety}
\label{sec:ptp}

\subsection{Ordering, matching, contracts}

\ampi{} keeps MPI's non-overtaking guarantee: two messages from the same rank to
the same rank on the same communicator that both match a receive are received in
send order. The reference implementation orders by a monotone message id, giving
the stronger global-FIFO property, which is worth having because agent harnesses
are debugged by reading traces.

A \emph{contract} is a typed description with a structural part (required keys,
non-empty keys, token bounds, regular expressions) and a semantic part (a
natural-language postcondition carried to the consumer). Matching follows MPI: a
send's and a receive's contracts match iff their names and kinds agree and the
receiver's required set is a subset of the sender's; only those fields participate
in matching, so a receiver may legally under-specify.

Separating \emph{required} (key present) from \emph{non-empty} (key present and
carries a value) sounds pedantic and is not. In an early version we conflated
them, and a term sheet for a book section containing no proper nouns --- a correct
answer --- was rejected, making the rank retry, exhaust its retries, raise, and
abandon peers already blocked in a collective. One over-strict contract field
produced a whole-population hang.

\subsection{Views: the datatype idea that matters}

A \emph{view} names a projection of an artifact without materialising it:
\code{contiguous}, \code{vector(offset,count,stride)}, \code{indexed},
\code{keys}, \code{jsonpath}, \code{head}, \code{tail}, \code{digest\_only}. This
is MPI's second datatype job --- \code{MPI\_Type\_vector} exists so a program can
send a boundary column without packing a copy --- and it is the job every agent
framework omits. \code{vector} expresses a block-cyclic decomposition (give rank
$r$ every $p$-th chapter) with no coordinator materialising the whole.
\code{head}/\code{tail} have no MPI analogue; they are the escape hatch that makes
budget compliance always achievable, and because a silent truncation is a
correctness hazard the runtime traces when they fire.

\subsection{Transport modes and eager credit}

\begin{table}[t]
\centering\small
\begin{tabular}{lll}
\toprule
mode & payload delivered & charged to context \\
\midrule
\code{EAGER} & yes & yes \\
\code{RENDEZVOUS} & envelope only & no \\
\code{SYNCHRONOUS} & yes, after match & yes \\
\code{AUTO} & by eager limit & as chosen \\
\bottomrule
\end{tabular}
\caption{Transport modes. MPI hides this choice because both modes deliver the
same bytes; here they differ in whether the receiving agent's context is consumed,
so the choice is semantic.}
\label{tab:modes}
\end{table}

A rendezvous envelope carries the digest, token count, contract and a one-line
\emph{synopsis}. The synopsis is not decoration: without it the receiver must
fetch in order to decide whether to fetch.

Every rank publishes an \emph{unexpected-message budget}: a bound in tokens on the
volume of unmatched eager messages it accepts. A sender that would exceed the
destination's budget blocks, and raises \code{ERR\_CONTEXT\_OVERFLOW} if it cannot
proceed by its deadline. Both the stall and its resolution are traced.

\subsection{Context safety}

\begin{definition}
A program is \emph{context-safe} if it completes for every assignment of
unexpected-message budgets, however small.
\end{definition}

\begin{proposition}
A program in which every rank sends before any rank receives is context-safe iff
its payloads travel \code{RENDEZVOUS}.
\end{proposition}

This is a direct transplant of MPI's safe-program discipline, and the agent version
is sharper. In MPI, exceeding the eager pool deadlocks or exhausts memory; here the
scarce buffer is the receiving agent's context, and exceeding it silently degrades
the agent's output. Making the budget explicit and blocking on it converts an
invisible quality failure into a reported, attributable stall --- which is the most
useful thing a protocol can do about it, because the harness can then narrow a view
and retry. \cref{tab:transport} measures both halves.

\section{Collectives and the algebra of semantic reduction}
\label{sec:coll}

Every collective is expressed over point-to-point, so its message count, rounds and
fold depth are observable. We deliberately do not short-circuit collectives through
the fabric as a shared variable: the value of the abstraction is that the
\emph{shape} of the communication determines cost and quality, and a centralised
implementation would hide exactly those effects. The one exception is the
arrival-counting barrier, which is centralised precisely because it must be able to
\emph{name} the ranks that did not arrive.

\subsection{Barriers must have a policy}

An unconditional barrier is a liveness bug waiting to happen: the probability that
all $p$ agent ranks arrive within a fixed window falls with $p$, and ``the agents
are waiting for each other'' is the most frequently reported multi-agent
pathology~\cite{mast}. Every \ampi{} barrier carries a deadline and one of five
policies: \code{WAIT} (MPI semantics; debugging only), \code{RAISE},
\code{PROCEED} (continue with those who arrived, report absentees), \code{SHRINK},
\code{REVOKE}. The policy is the harness author's declaration of what a missing
participant \emph{means}: a missing chapter degrades a translation, a missing
module kills a build.

\subsection{Operators declare their algebra}

An operator declares commutativity, associativity in
$\{\textsc{exact},\textsc{approx},\textsc{none}\}$, idempotence, and an identity.
The runtime enforces:

\begin{itemize}\itemsep2pt
\item \textsc{none} permits only the serial chain; requesting a tree is an error,
not a silent quality regression.
\item A non-commutative operator is evaluated in rank order. A binomial tree
satisfies this only when the root is rank 0; for any other root the runtime
\emph{refuses} the tree rather than silently evaluating in rotated order.
\item The serial left fold is the reference semantics: for an \textsc{exact}
operator every algorithm must reproduce it, and our test suite checks this across
nine population sizes and every algorithm.
\item Every reduction reports the \emph{fold depth} it induced.
\end{itemize}

\subsection{Fidelity decays in depth, not in width}

\begin{table}[t]
\centering\small
\begin{tabular}{lccc}
\toprule
algorithm & rounds & applications & fold depth \\
\midrule
\code{chain}    & $p-1$                  & $p-1$                 & $p-1$ \\
\code{flat}     & $1$                    & $p-1$ (all at root)   & $p-1$ \\
\code{binomial} & $\lceil\log_2 p\rceil$ & $p-1$                 & $\lceil\log_2 p\rceil$ \\
\code{kary}($k$) & $\lceil\log_k p\rceil$ & $p-1$ transfers, $\lceil\frac{p-1}{k-1}\rceil$ applications & $\lceil\log_k p\rceil$ \\
\bottomrule
\end{tabular}
\caption{Reduce algorithms. The first three perform the same number of operator
applications and differ only in rounds and fold depth: if quality were governed by
the number of applications they would score alike, and if by depth the tree wins.
\code{kary} additionally reduces the number of applications, because a variadic
operator combines all $k$ children in one call.}
\label{tab:reduce-alg}
\end{table}

\begin{table}[t]
\centering\small
\begin{tabular}{rrrrr}
\toprule
fan-in $k$ & rounds & messages & fold depth & modelled $F(d)$ \\
\midrule
2  & 6 & 63 & 6 & 0.735 \\
4  & 3 & 63 & 3 & 0.857 \\
8  & 2 & 63 & 2 & 0.902 \\
16 & 2 & 63 & 2 & 0.902 \\
\bottomrule
\end{tabular}
\caption{$k$-ary reduction at $p=64$, $\delta=0.05$. The message count does not
depend on $k$ --- every non-root rank sends exactly once --- so widening the tree is
free in volume and pays in both latency and fidelity. This is the opposite of the
trade MPI faces, and it follows from agent ranks not being single-ported.}
\label{tab:fanin}
\end{table}

MPI requires associativity and then reorders freely; floating-point \code{MPI\_SUM}
is consequently not bitwise reproducible, and this is tolerated because the
discrepancy is bounded by rounding error~\cite{Demmel15}. An operator implemented
by a language model is worse in four ways: it is not associative, it is lossy, it is
expensive, and it is non-deterministic. Crucially its loss compounds with the
\emph{depth} of the reduction, because depth counts how many times an item is
re-summarised on the way to the root. Modelling surviving fidelity of a $d$-deep
fold as $F(d)=(1-\delta)^d$ for a measured per-application loss $\delta$,
\cref{tab:reduce-alg} predicts that the tree dominates.

This inverts the usual expectation that parallel aggregation trades quality for
speed. It holds only because of the immutability decision in
\cref{sec:model}: were interior ranks relaying regenerated text, a tree broadcast
would incur its own depth-proportional drift, and the fidelity ordering would
reverse. \Cref{tab:fidelity} measures retention.

\paragraph{Binary is the wrong arity.} The reasoning above stops one step short.
Every $\log_2 p$ factor in MPI's collective layer descends from a single hardware
premise: a process is \emph{single-ported}, so two incoming messages cost two
transfers however the tree is shaped, and a wider node is pure loss. Agent ranks
violate that premise completely. A prompt can carry $k$ artifacts as easily as two,
and an operator with a \emph{variadic} kernel combines all $k$ in one application.
The constraint that replaces port count is context capacity, so the optimal fan-in
is $k^{*}=\lfloor C(1-h)/s \rfloor$ for a budget $C$, per-input size $s$, and a
reserve fraction $h$, and the fold depth becomes $\lceil\log_{k^{*}} p\rceil$.

\Cref{tab:fanin} gives the consequence. Widening from $k{=}2$ to $k{=}8$ at $p=64$
cuts rounds and fold depth from six to two while sending exactly the same 63
messages, because every non-root rank sends once regardless of arity. In MPI,
widening a reduction tree trades latency against bandwidth; here it costs nothing
and buys both latency and fidelity. \ampi{} therefore ships \code{kary} with the
fan-in derived from the declared context budget, and \code{optimal\_fanin} caps it,
because the relationship is not monotone in practice --- a model handed fifty
documents in one call attends to them unevenly, so reduced depth is eventually
offset by within-call neglect. Where that cap belongs is an empirical question and
the benchmark measures it.

\paragraph{Wide trees need variadic operators, and the runtime says so.} A $k$-ary
tree given only a binary kernel degrades to a local left fold at each node, which
spends exactly the $k-1$ successive applications the wide tree was meant to avoid.
\ampi{} records which path was taken, because the failure is silent otherwise: the
message count and round count look correct while the fold depth quietly reverts to
$p-1$.

\paragraph{A third effect has no MPI counterpart.} A summarising operator given a
fixed output budget compresses two inputs and sixteen inputs to the same length, so
the root's answer over-represents whichever subtree merged last --- an unfairness
\code{MPI\_SUM} cannot exhibit, since it weights every contribution equally by
construction. \ampi{} therefore exposes the number of leaves an accumulator
represents so a budget can be allocated proportionally, and reports per-rank
retention so the unfairness is visible rather than inferred.

\subsection{Prefer exact operators}

\code{UNION} --- set union, or key-wise union with deterministic conflict retention
--- is exact, associative, commutative and idempotent. A glossary merged with
\code{UNION} is byte-identical at every rank regardless of tree shape; a glossary
merged by an agent is not. Harnesses should prefer exact operators wherever the
combination is genuinely set-like, exactly as an MPI programmer prefers
\code{MPI\_SUM} to a user-defined operator. Idempotence additionally matters
because a rank may be re-executed after a suspected but unconfirmed failure, and a
non-idempotent merge applied twice double-counts.

\subsection{The allreduce divergence hazard}

\code{reduce\_bcast} computes one result and broadcasts it; \code{recursive\_doubling}
has each rank compute the result itself, which is latency-optimal in MPI. The
difference matters here, but not for the reason one would first assume, and the
distinction is worth drawing precisely because it separates a hazard that can be
engineered away from one that cannot.

The obvious worry is that under an independent-fold algorithm each rank folds in its
own \emph{arrival order}, so a non-associative operator yields $p$ different answers.
That worry is real and it is removable: ordering the operands of every pairwise
combine by \emph{rank} rather than by arrival makes both partners evaluate the
identical expression tree, at no cost. It is the same reasoning MPI uses to permit
trees for operators declared non-commutative, and with it a
\code{recursive\_doubling} allreduce over a deterministic operator --- however lossy
or order-sensitive --- gives every rank the same result. We verify this, because the
property is invisible to any test that uses an exact operator.

What ordering cannot fix is that a semantic operator is \emph{nondeterministic}:
every rank evaluates the same expression and gets different text. So the population
still diverges under \code{recursive\_doubling}, and \textbf{silently disagrees
about the value it just agreed on} --- but the cause is the operator, not the
algorithm. The only remedy is for exactly one rank to evaluate and the rest to
receive the result by handle, which is what \code{reduce\_bcast} does. \ampi{}
therefore defaults lossy operators to \code{reduce\_bcast}, paying a factor of two
in rounds to buy the property that agreement means agreement, and \emph{flags}
rather than forbids the alternative, since a harness whose operator happens to be
deterministic should still be allowed the faster algorithm.

\subsection{Scan, and topologies}

Scan is the most under-used collective for agent work. A task with sequential
dependence but parallel bulk --- translate chapter $i$ consistently with names
established in $0..i{-}1$; write a section that must not contradict earlier ones ---
\emph{is} a prefix computation. Running it sequentially costs $p$ agent latencies;
ignoring the dependence produces inconsistent output. Hillis--Steele recursive
doubling delivers every prefix in $\lceil\log_2 p\rceil$ rounds at the price of
$p\log p$ operator applications: it trades \emph{more} operator work for
\emph{fewer} rounds, the opposite of the usual trade, and the right direction
whenever per-invocation latency dominates (\cref{sec:cost}).

A virtual topology restricts who may communicate with whom, and that restriction is
the main lever on the quadratic blow-up that kills multi-agent systems. A full
conversation over $p$ agents costs $\Theta(p^2)$ messages and $\Theta(p^2 n)$
tokens; over a ring or a degree-2 review graph the same information costs
$\Theta(p)$ and $\Theta(pn)$. Our test suite measures the gap directly: at $p=8$ a
pairwise alltoall sends 56 messages where a degree-2 neighbourhood allgather sends
16. Halo exchange is the most valuable non-obvious transfer from HPC: a task whose
parts each depend on their neighbours \emph{looks} sequential, and expressed as a
stencil it is one barrier-separated parallel step plus a constant-degree boundary
exchange.

\section{One-sided state: the shared blackboard}
\label{sec:rma}

Every multi-agent system grows a shared scratchpad --- a design document, a
glossary, an interface file, a task board --- and every one gets the concurrency
wrong the same way: two agents read it, both edit a private copy, and the second
write discards the first's work. MPI-3's one-sided chapter supplies the
vocabulary~\cite{hoefler15-rma,dinan16-rma}.

A \emph{window} is a named, explicitly created region over a communicator, divided
into \emph{slots}; slots rather than byte offsets are the unit of locking and
versioning, because the natural granularity of agent-shared state is ``the section
about error handling'', not ``bytes 4096--8191''. The operations that matter are
distinguished precisely:

\begin{itemize}\itemsep2pt
\item \code{put} --- blind overwrite: the operation that loses work. With
\code{expect\_version} it becomes a compare-and-swap.
\item \code{accumulate(op)} --- atomic read-modify-write, so concurrent
contributions combine instead of clobbering. \emph{No lock is needed.} The operator
must be exact: the read-modify-write happens inside an atomic section and an
implementation cannot hold one across a model call.
\item \code{compare\_and\_swap}, \code{fetch\_and\_op} --- optimistic concurrency
and atomic counters; the latter is how a harness implements dynamic
self-scheduling, which matters more for agents than for CPUs because per-item cost
varies by an order of magnitude.
\item \code{lock}/\code{unlock} with a \emph{lease}, \code{fence} (collective epoch
boundary), \code{sync} (discard private copies).
\end{itemize}

\paragraph{\code{SEPARATE} is the right default.} MPI-3 distinguishes a
\emph{unified} model, in which a rank always sees the latest value, from a
\emph{separate} model, in which each rank has a private copy reconciled only at
synchronisation points. For an agent the private copy is not a hardware artefact:
it is the copy of the document sitting in the agent's context from ten minutes ago,
and it \emph{is} stale, because peers have edited it since. \code{sync} is
``re-read before you edit''. \ampi{} therefore defaults to \code{SEPARATE}, tracks
each rank's observed version per slot, and records a \textbf{staleness violation}
when a rank writes from a version that is no longer current.

Recording the violation is more valuable than preventing it. A harness that
legitimately overwrites must still be able to, so refusing the write would break
correct programs; staying silent would reproduce the lost-update bug. Our test
suite asserts the \emph{failure}: eight ranks that read, modify and blindly write a
shared document produce a final value with fewer than eight contributions, and the
runtime counts the stale writes that caused it. The same test with
\code{accumulate} loses nothing and records zero stale writes.

Leases have no MPI counterpart because a dead MPI process takes the job with it.
Here, a lock held by a vanished agent would deadlock the window forever, so every
lock is leased and reclamation is traced --- a held-then-stolen lock is exactly the
situation in which two agents believe they own the same file.

\paragraph{Semantic updates serialise, and the trace says so.} A harness needing
judgement in a combine must \code{lock}, \code{get}, reason, \code{put} with
\code{expect\_version}, \code{unlock}. Making that explicit is the point: a harness
that puts a semantic critical section on its critical path has built a sequential
program with extra steps, and the trace shows it as lock wait time.

\section{Failure model and fault tolerance}
\label{sec:ft}

\begin{table*}[t]
\centering\small
\begin{tabular}{llll}
\toprule
class & description & detector & mitigation \\
\midrule
F1 fail-stop & crashed, killed, past a hard deadline & lease expiry & re-incarnate, or shrink \\
F2 fail-slow & alive but past its latency budget & deadline from the run's own distribution & hedge with a duplicate \\
F3 fail-noisy & violates its structural contract & contract check & retry with the diagnosis \\
F4 fail-plausible & well-formed, confident, wrong & \textbf{independent verification only} & replicate and compare; oracle \\
F5 fail-greedy & exhausted its context budget & context accounting & rendezvous transport plus views \\
F6 adversarial & deliberate protocol deviation & --- & out of scope \\
\bottomrule
\end{tabular}
\caption{The \ampi{} failure model. MPI's fault-tolerance work assumes F1;
importing that assumption is why much agent-reliability engineering consists of
retry loops that cannot help. F4 is the load-bearing row.}
\label{tab:failures}
\end{table*}

\subsection{Six classes, and why F4 dominates}

\Cref{tab:failures} gives the model. The row that matters is F4. A rank returning a
confident wrong answer is undetectable by any amount of timeout tuning, retry logic
or schema validation, because none of those examine whether the answer is
\emph{right}. It is the exact analogue of silent data corruption, and HPC's answer
to silent data corruption is not checkpoint/restart --- which preserves the
corruption faithfully --- but algorithm-based fault tolerance: carry redundant
information that lets the computation check
itself~\cite{huang84abft}. Consequently \ampi{} makes verification a first-class
concept, and we argue the quantity to optimise is the \emph{verification budget},
not the checkpoint interval that Daly's formula
optimises~\cite{daly06checkpoint}.

Detection is by lease, and is \emph{eventually perfect} rather than perfect: a slow
rank is indistinguishable from a dead one. FLP says consensus is unattainable with
an unreliable detector in an asynchronous system~\cite{fischer85flp}, and we take
the same escape as HPC --- assume partial synchrony with a known lease bound, accept
that an expired lease \emph{declares} a rank dead whether or not it is, and use
\code{revoke} to make the declaration self-fulfilling so the population's view
stays consistent even when the detector was wrong~\cite{chandra96failuredetectors}.
Straggler thresholds are derived from the run's own latency distribution, because
the median-to-tail spread of an agent invocation is routinely an order of magnitude
and is driven by output length rather than queueing~\cite{dean13tail}; a fixed
timeout is either useless or harmful.

\subsection{Revoke, shrink, agree}

\code{revoke} makes a communicator permanently unusable for \emph{every} member;
operations already blocked inside a collective raise \code{ERR\_REVOKED}. This is
ULFM's least obvious and most necessary primitive. Without it, a rank that notices
a failure cannot rescue the others, because they are all blocked waiting for the
dead peer --- and it cannot first reach agreement, because the group is broken.
Revocation is irreversible: a communicator that could be un-revoked would let two
ranks disagree about whether it is usable. Implementing this correctly required a
subtlety worth recording --- a blocked call must consult shared state rather than a
cached flag, or it never learns it has been revoked.

\code{shrink} builds a communicator over the survivors, renumbered contiguously.
\code{shrink\_in\_place} instead marks ranks inactive and keeps the numbering ---
FT-MPI's \code{BLANK} mode rather than its \code{SHRINK} mode~\cite{fagg00-ftmpi}
--- and is often the better trade here, because renumbering invalidates a
rank-to-work mapping that is expensive to recompute.

\code{agree} returns the conjunction over live ranks, consistently: either every
survivor returns the same value or all raise. This is what makes ``did the
integration step succeed?'' answerable when some builders died mid-step; an
allreduce over logical AND would hang.

\subsection{Redundancy and supervision}

\code{replicate\_and\_compare} reports the modal answer, the number of distinct
answers, and the consensus fraction, comparing on a harness-supplied projection
because byte equality is meaningless for prose. One caution: replicating a language
model does not give independent failures, since correlated errors are the norm.
Agreement is evidence, not proof, which is why the consensus fraction is reported
rather than a bare boolean.

MPI contributes the communication algebra and has no recovery policy; Erlang/OTP
contributes the recovery discipline and no collectives~\cite{armstrong03erlang}. The
two are orthogonal --- each field solved what the other ignored --- and \ampi{}
takes both, implementing OTP's \code{one\_for\_one}, \code{one\_for\_all} and
\code{rest\_for\_one} strategies bounded by a restart intensity. The bound matters
especially here: an unbounded restart loop converts a local fault into a cost
overrun.

\subsection{The rule harness authors break}

\begin{quote}
\emph{A local failure must not remove a rank from a collective.}
\end{quote}

A rank that cannot compute its contribution must still enter the collective,
contributing a degraded value or an identity element, and record the degradation.
This is the rule most easily violated and most expensive to violate: if a local
exception propagates out of \code{rank\_main}, that rank never reaches the
collective its peers are already blocked inside, and one recoverable local failure
becomes a whole-population hang. We hit exactly this in our own translation harness
before adding the \code{safe\_agent} wrapper. Escalation must be a deliberate
decision made by a barrier policy or a supervisor, never an accident of exception
propagation.

\section{Cost model}
\label{sec:cost}

Volume is measured in tokens; a message of $n$ tokens costs
\[ T(n)=\alpha+n\beta, \qquad C(n)=n\gamma . \]
Three properties distinguish this from Hockney's model.

\paragraph{$\alpha$ is enormous, so rounds matter more than volume.} For MPI,
$\alpha$ is microseconds and bandwidth gigabytes per second, so the ratio rewards
reducing message volume. For agent ranks we measure $\alpha_{p50}$ in the tens of
seconds and effective bandwidth in the tens of tokens per second, so the ratio
rewards reducing \emph{rounds}, and latency-optimal algorithms win over a wider
range of sizes than in MPI. \Cref{tab:calibration} gives the measured constants;
the crossover $\alpha/\beta$ lands at a few hundred tokens --- the size of a typical
agent artifact --- so neither term may be neglected.

\paragraph{Price is an independent axis.} Wall time and money are not proportional:
running $p$ ranks divides time by up to $p$ and divides price by nothing, and
coordination messages are pure added price. A harness has two objectives, and we
report both everywhere.

\paragraph{Fidelity is a third axis.} $F(d)=(1-\delta)^d$, as in \cref{sec:coll}.

We report Amdahl speedup, where the serial fraction $f$ is dominated by phases that
must be done by one rank --- plan, integrate, arbitrate --- and those do not
shrink~\cite{amdahl67}; the Universal Scalability Law
$S(p)=p/(1+\sigma(p-1)+\kappa p(p-1))$, whose \emph{coherency} term $\kappa$ is
positive by construction for an all-to-all conversation, so its throughput has a
maximum and then declines~\cite{gunther07usl}; and Karp--Flatt, which reveals a
growing coordination cost that a plain speedup plot hides~\cite{karp90metric}.

\section{Implementation}
\label{sec:impl}

The reference implementation is 6\,900 lines of Python 3.11 with no third-party
dependencies. Layers: a durable fabric (SQLite in WAL mode with
\code{BEGIN IMMEDIATE} writers and jittered retry, plus a content-addressed blob
store with fsync-and-rename); a transport with matching, non-overtaking sequence
numbers and eager credit; communicators; twelve collectives with 27 algorithms;
windows; fault tolerance; a cost model and calibrated simulator; and a CLI so an
external agent can act as a rank.

Executors are pluggable, mirroring MPI's separation of interface from process
manager: a deterministic function, a calibrated simulator, a recorded replay, and a
\emph{broker} that publishes invocations for external workers to claim. The broker
is a \emph{pull} queue, because a pushed invocation would require the harness to
know how to start an agent and would couple it to a vendor; pulling lets the
population be launched, scaled and re-incarnated entirely outside the harness,
which is required because an agent session's lifetime is controlled by its host.
The queue is per rank, because a rank is a durable role with accumulated state.

\paragraph{Tracing is unconditional} and in the same transaction as the state change
it describes. MPI's tooling interfaces are opt-in and out-of-band, which is
reasonable when a run can be repeated cheaply. An agent run cannot: it is expensive
and not reproducible, so a bug that was not traced cannot be investigated.

\paragraph{Bugs the design caught.} Building the cost model and the implementation
independently and cross-checking them found four real defects: a binomial scatter
whose parent was computed from the highest rather than the lowest set bit (so it
disagreed with the broadcast tree and deadlocked for $p\geq4$); a chain reduce that
folded from $\text{root}+1$ and so broke rank order for non-commutative operators;
an allreduce that reported the broadcast's fold depth instead of the reduction's;
and two cost formulas that over-counted messages at non-power-of-two sizes by up to
25\%. An implementation whose measured counts disagree with its own model has a bug
in one of the two, and having both is how the discrepancy gets found.

\section{Evaluation}
\label{sec:eval}

We evaluate four questions. \textbf{Q1}: does the implementation match its cost
model? \textbf{Q2}: do the protocol mechanisms behave as specified under context
pressure and failure? \textbf{Q3}: does an \ampi{} harness deliver a useful
parallel speedup on a weakly coupled task, and do its coordination mechanisms
matter? \textbf{Q4}: does it help on a strongly coupled task, where naive
parallelism fails?

All agent-executed runs use a population of real agent workers claiming invocations
from the broker; the worker bootstrap prompt is in the artifact and mentions
neither collectives nor the experiment (\cref{sec:model}). Protocol-level results
(Q1, Q2) use deterministic executors so that they are free and regression-testable.
Every quality number is computed by code from the agents' artifacts, never from a
model's opinion of its own output; model-judged scores would make a comparison
between protocol variants unreplicable.

\subsection{Q1: model validation}
\label{sec:eval-micro}

\begin{table}[t]
\centering\scriptsize
\begin{tabular}{lrrrrr}
\toprule
collective & algorithms & sizes & configs & message count & fold depth \\
\midrule
allgather & 3 & 13 & 31 & 31/31 & 31/31 \\
allreduce & 2 & 13 & 26 & 26/26 & 26/26 \\
alltoall & 3 & 13 & 39 & 39/39 & 39/39 \\
barrier & 2 & 13 & 26 & 26/26 & 26/26 \\
bcast & 3 & 13 & 39 & 39/39 & 39/39 \\
gather & 2 & 13 & 26 & 26/26 & 26/26 \\
reduce & 3 & 13 & 39 & 39/39 & 39/39 \\
scan & 2 & 13 & 26 & 26/26 & 26/26 \\
scatter & 2 & 13 & 26 & 26/26 & 26/26 \\
\midrule
\textbf{total} & & & 278 & 278/278 & 278/278 \\
\bottomrule
\end{tabular}

\caption{Measured versus closed-form message counts and fold depths, over every
implemented (collective, algorithm) pair at population sizes 2--16.}
\label{tab:model}
\end{table}

\begin{table}[t]
\centering\scriptsize
\begin{tabular}{lrl}
\toprule
parameter & value & meaning \\
\midrule
$\alpha_{p50}$ & 59.5 s & median agent invocation latency \\
$\alpha_{p99}$ & 389.6 s & tail latency \\
$\beta^{-1}$ & 13.1 tok/s & marginal output rate \\
$\alpha/\beta$ & 782 tok & latency/volume crossover \\
fabric & 0.0030 s & per-operation protocol cost \\
$\alpha/\text{fabric}$ & 19,840$\times$ & agent cost over protocol cost \\
\bottomrule
\end{tabular}

\caption{Cost-model constants calibrated from the translation runs. The last row is
the ratio of an agent invocation to a protocol operation, and it is why keeping
protocol state in a durable external store is free.}
\label{tab:calibration}
\end{table}

\Cref{tab:model} compares measured message counts and fold depths against the
closed-form formulas. \Cref{tab:calibration} gives the calibrated constants. The
$\alpha$-to-fabric ratio is the design's licence: the protocol's per-operation cost
is negligible against an agent invocation, so durability, tracing and contract
checking can all be unconditional.

\subsection{Q2: context safety and failure recovery}

\begin{table}[t]
\centering\scriptsize
\begin{tabular}{rccrccr}
\toprule
& \multicolumn{3}{c}{eager} & \multicolumn{3}{c}{rendezvous} \\
\cmidrule(lr){2-4}\cmidrule(lr){5-7}
payload (tok) & completes & error & max ctx & completes & max ctx & tok deferred \\
\midrule
128 & \checkmark & -- & 168 & \checkmark & 0 & 2,688 \\
512 & \checkmark & -- & 674 & \checkmark & 0 & 10,784 \\
2,048 & \checkmark & -- & 2,695 & \checkmark & 0 & 43,120 \\
8,192 & $\times$ & ERR\_CONTEXT\_OVERFLOW & 0 & \checkmark & 0 & 172,464 \\
32,768 & $\times$ & ERR\_CONTEXT\_OVERFLOW & 0 & \checkmark & 0 & 689,856 \\
\bottomrule
\end{tabular}

\caption{Ring exchange in which every rank sends before any rank receives, at
$p=8$ with an unexpected-message budget of \NEagerLimit{}
tokens. Eager transport completes only while the payload fits the budget; rendezvous
completes at every size.}
\label{tab:transport}
\end{table}

\begin{table}[t]
\centering\scriptsize
\begin{tabular}{lccccc}
\toprule
policy & survivors complete & detect (s) & absentees named & shrink (s) & agreement consistent \\
\midrule
\texttt{wait} & $\times$ & 15.25 & $\times$ & -- & \checkmark \\
\texttt{raise} & $\times$ & 15.25 & $\times$ & -- & \checkmark \\
\texttt{proceed} & \checkmark & 15.20 & \checkmark & -- & \checkmark \\
\texttt{shrink} & \checkmark & 15.20 & \checkmark & 0.000 & \checkmark \\
\bottomrule
\end{tabular}

\caption{One rank of eight fails silently before a barrier. \code{WAIT} is MPI's
only option and never completes.}
\label{tab:faults}
\end{table}

\Cref{tab:transport} confirms the context-safety proposition with a sharp
threshold, and shows the failure is \emph{reported} as
\code{ERR\_CONTEXT\_OVERFLOW} rather than manifesting as a hang or a silent
truncation. \Cref{tab:faults} shows that \code{WAIT} and \code{RAISE} do not let the
survivors complete, while \code{PROCEED} and \code{SHRINK} do, correctly name the
absentee, and reach consistent agreement.

\subsection{Q3: parallel book translation}

\begin{table*}[t]
\centering\scriptsize
\begin{tabular}{lrrrrrrrrr}
\toprule
configuration & consist. & diverg. & verified & para. & CJK & wall (s) & calls & tokens & USD \\
\midrule
glossary, halo & 1.000 & 0/0 & 0.000 & 0.000 & 0.000 & 7,200 & 10 & 13,925 & 0.073 \\
\bottomrule
\end{tabular}

\caption{Translation ablations at the largest population, with real agent ranks.
``consist.'' is the fraction of entities rendered identically by every section that
rendered them; ``verified'' is the fraction of reported renderings that actually
occur in the reporting section's text, which is a direct measurement of
fail-plausible behaviour.}
\label{tab:translation}
\end{table*}

\begin{table}[t]
\centering\scriptsize
\begin{tabular}{rrrrrrrrr}
\toprule
$p$ & wall (s) & speedup & efficiency & Karp--Flatt & calls & tokens & USD & consist. \\
\midrule
1 & 577 & 1.00 & 1.00 & 0.000 & 16 & 33,756 & 0.161 & 1.000 \\
2 & 410 & 1.41 & 0.70 & 0.422 & 18 & 38,253 & 0.186 & 1.000 \\
4 & 308 & 1.87 & 0.47 & 0.380 & 22 & 46,926 & 0.231 & 1.000 \\
8 & 167 & 3.46 & 0.43 & 0.188 & 24 & 52,084 & 0.260 & 1.000 \\
\bottomrule
\end{tabular}

\caption{Strong scaling: identical total work, decreasing population. USL fit:
\NUslFit.}
\label{tab:trscaling}
\end{table}

Translating a book looks embarrassingly parallel and is not: sections disagree about
proper nouns, drift in register, and show visible seams. The dependence is weak but
real, and has the shape of a stencil --- mostly local work, a little that must be
globally consistent, a little that must agree with immediate neighbours. Our harness
is the corresponding HPC decomposition (\cref{sec:model}): \code{scatterv},
agent-local term extraction, \code{allreduce} with \code{UNION} to agree a glossary
in $2\lceil\log_2 p\rceil$ rounds with no coordinator reading any section,
translation, \code{halo\_exchange} on a 1-D Cartesian topology to reconcile seams,
and \code{gather} by handle so the root never admits $pn$ tokens.

\paragraph{The protocol delivers a real speedup with no coherency ceiling.}
\Cref{tab:trscaling} shows $3.46\times$ speedup at $p=8$ on identical total work,
at 43\% efficiency. The interesting number is the USL fit, \NUslFit: the
\emph{contention} coefficient $\sigma$ is substantial at $0.22$, and the
\emph{coherency} coefficient $\kappa$ is zero. That is the distinction the model
exists to draw. A positive $\kappa$ would mean the coordination cost grows
superlinearly and throughput has a maximum beyond which adding agents makes the
system slower --- the signature of an all-to-all conversation, and the ceiling that
group-chat architectures hit. Measuring $\kappa=0$ says this harness's cost is a
constant serial fraction, not a ceiling, which is precisely what restricting
communication to collectives and a bounded-degree topology is supposed to buy. The
Karp--Flatt metric agrees: it does not grow with $p$.

\paragraph{The coordination machinery has a measured price.} Reading
\cref{tab:translation} as a cost table rather than a quality table: relative to
translating with neither mechanism, adding the glossary \code{allreduce} and the
halo exchange costs \NCoordTime\% more wall time and \NCoordTokens\% more tokens.
That is what agreement costs, quantified, and it is the number a harness author
needs before deciding to pay it.

\paragraph{Exact operators are worth reaching for, and the margin is large.} The
same glossary agreement performed by a semantic \code{allreduce} --- an
LLM-implemented merge --- rather than by the exact, idempotent \code{UNION} cost
$\NSemTime\times$ the wall time, $\NSemTokens\times$ the tokens and
$\NSemPrice\times$ the price, and produced glossaries of the same size with
identical measured quality. It also raised the agent-call count from 24 to
\NSemCalls, because a semantic reduction adds one model invocation per tree node on
top of the per-section work. This is the strongest practical result in the
experiment and it is a direct analogue of preferring \code{MPI\_SUM} to a
user-defined operator: when a combination is genuinely set-like, an exact operator is
not merely cheaper but exactly reproducible, whereas the semantic version is neither.
Choosing \code{recursive\_doubling} over \code{reduce\_bcast} for the exact operator,
by contrast, changed almost nothing (158\,s versus 167\,s), which is what the cost
model predicts when $p$ is small enough that a factor of two in rounds is within
run-to-run variance.

\paragraph{An honest negative result.} On this text, they did not have to. All
configurations achieved perfect entity consistency, including the one with no
glossary exchange at all, because the pre-registered entities are famous enough
that competent models render them identically without coordination. We report this
rather than substituting a metric after seeing the data. Two things follow, and
both are useful. First, our consistency metric is saturated on canonical entities
and a sensitive replication must score rare names
(\cref{sec:limitations}). Second, and more generally: \emph{a coordination
mechanism only pays where the participants would otherwise disagree}, so a harness
should spend agreement on the terms that are genuinely ambiguous and not on the
ones everybody already knows. The protocol makes that a decision the author can
price, which is more than a framework that broadcasts everything to everyone can
offer.

Two quality checks did not saturate and are worth recording. Every reported entity
rendering actually occurred in the text that reported it, so on this task we
observed no fail-plausible behaviour on the one axis where we could
mechanically detect it. And paragraph structure was preserved exactly in every
section, which is the contract the harness enforced and which a free-running
translation agent routinely violates.

\subsection{Q4: collaborative software development}

\begin{table*}[t]
\centering\scriptsize
\begin{tabular}{lrcrrrrrrr}
\toprule
configuration & $p$ & passed & lines & defensive & /kloc & wall (s) & calls & tokens & USD \\
\midrule
shared interfaces & 8 & 60/60 & 1,954 & 7 & 3.58 & 499 & 32\rlap{$^\ddagger$} & 161,989 & 1.145 \\
\textbf{no shared interfaces} & 8 & 59/60 & 2,693 & 33 & 12.25 & 529 & 32\rlap{$^\ddagger$} & 167,174 & 1.377 \\
single agent ($p{=}1$) & 1 & 60/60 & 1,752 & 1 & 0.57 & 1,213 & 2 & 31,945 & 0.438 \\
\midrule
shared interfaces, \emph{vague spec} & 8 & 60/60 & 2,203 & 2 & 0.91 & 402 & 16 & 86,473 & 0.647 \\
\textbf{no shared interfaces}, \emph{vague spec} & 8 & 60/60 & 3,564 & 15 & 4.21 & 848 & 32 & 191,602 & 1.643 \\
\bottomrule
\end{tabular}

{\footnotesize $^\ddagger$ This run's first round already passed the acceptance suite, but the suite's report failed to parse and was recorded as an import failure, so \code{agree} withheld the build and a second round was spent repairing a defect that did not exist. The quality columns come from an offline re-evaluation against one oracle and are unaffected; the wall, call, token, and price columns include that extra round, so they are not comparable with single-round rows.}

\caption{Collaborative implementation of the \code{minidb} SQL engine against a
fixed 59-case acceptance suite written before the run. ``per round'' is the number
of cases passing after each integration round.}
\label{tab:software}
\end{table*}

\begin{table}[t]
\centering\scriptsize
\begin{tabular}{lrrrrrrrr}
\toprule
configuration & locks & contended & wait (s) & max wait & puts & stale & stale rate & accum. \\
\midrule
shared interfaces & 0 & 0 & 0 & 0.00 & 8 & 0 & 0.000 & 0 \\
\bottomrule
\end{tabular}

\caption{Contention on the shared interface window.}
\label{tab:contention}
\end{table}

The complement to translation. Eight agents each handed one module of a
specification produce eight modules that do not compose, because each invented its
neighbours' interfaces. The harness is organised around exactly that failure:
\code{bcast} the specification; publish interfaces into a window slot per module and
\code{fence}; read dependencies' interfaces in a fresh access epoch; implement;
\code{barrier}; run the acceptance suite at rank 0; \code{scatter} per-module blame
(not broadcast everything --- a rank should receive its own failures);
\code{neighbor\_allgather} for degree-2 cross review; renegotiate an interface under
an exclusive lock; \code{agree} on whether the build is green.

The oracle is a 59-case acceptance suite written before the run and never shown to
the agents as a target to optimise, run out of process because generated code can
hang or exhaust memory. It is also our model of verification (\cref{sec:ft}): each
case is an independent check that a confident output actually behaves as specified.

\Cref{tab:software} reports pass rates and \cref{tab:contention} the window
contention, which bounds the achievable speedup Amdahl-style.

\paragraph{The parallel population buys latency with money, at equal quality.} The
headline comparison is against the baseline a multi-agent system has to beat. Eight
agents reached \NSwParPassed{} in \NSwParWall\,s using \NSwParCalls{} agent
invocations; one agent writing all eight modules itself reached \NSwSoloPassed{} in
\NSwSoloWall\,s using \NSwSoloCalls. That is a $\NSwSpeedup\times$ latency reduction
at \NSwEff\% parallel efficiency, bought for $\NSwPrice\times$ the price and
$\NSwTokens\times$ the tokens.

We think this is the right way to report a multi-agent result, and it is not the way
they are usually reported. The population did not achieve anything the single agent
could not; the task is specified precisely enough that one agent gets there in two
passes. What the population bought was wall-clock time, and what it paid was money and
tokens --- which is exactly what \cref{sec:cost}'s insistence on price as an
\emph{independent} axis predicts, since running $p$ ranks divides time by up to $p$ and
divides price by nothing. A harness author facing this trade needs both numbers, and a
paper that reported only the speedup would be concealing the more important one. It
also explains, without any appeal to model capability, the widely reported finding that
multi-agent architectures frequently fail to beat a single
agent~\cite{cognitiondont,lmteams26}: on a well-specified task with no latency
pressure, there is nothing for them to win.

\paragraph{The single agent could verify continuously; no rank could.} The baseline's
result comes with an asymmetry we should have anticipated and did not, and it is
structural rather than an artefact. The single agent, asked for all eight modules at
once, wrote them to a scratch directory, wrote itself a 49-case test suite covering the
specification's NULL logic, aggregate rules, ordering and error-wrapping guarantees, and
iterated until they passed --- \emph{before} submitting anything. It reported doing so
unprompted.

No rank in the parallel population could do that. A rank owns one module of eight and
cannot execute the system, so for it the first opportunity to observe whether anything
works is the integration barrier. Decomposition does not merely divide the work; it
divides the work into pieces that are individually \emph{untestable}, and it moves
verification from something continuously available to something available once per
round. That is a real disadvantage of parallel agent development, it is invisible in a
pass-rate table, and it is not fixed by better prompting.

It does suggest a concrete protocol-level remedy, which we did not implement and would
in a revision: verification should be reachable \emph{inside} a rank's loop rather than
only at the barrier. A rank could pull the current contents of the shared window,
assemble the system from its peers' latest published artifacts, and run the acceptance
suite locally --- speculatively, against possibly-stale peers, in the spirit of the
\code{SEPARATE} memory model where a rank knowingly reasons about a private copy. The
protocol already has the pieces; the harness simply did not use them, and
\cref{sec:ft}'s claim that verification is the load-bearing fault-tolerance mechanism
implies it should have.

\paragraph{Eight agents composed a working system on the first integration.} With
interfaces published into the RMA window, the population produced roughly 1\,950
lines across eight modules --- tokeniser, AST, parser, planner, expression and
aggregate functions, execution engine, public API --- and passed the entire
acceptance suite at the first integration barrier. No agent saw another agent's code
at any point; each was given the specification, its own module's assignment, and the
published interfaces of the modules it depends on. That a parser written against a
declared \code{tokenize} signature composes with a planner written against a declared
AST, without either author seeing the other's implementation, is the whole claim of
the interface-publication phase, and it is the outcome an ordinary "eight agents,
eight modules" prompt does not produce.

\paragraph{Withholding interfaces does not cost correctness; it costs defensive
coupling.} This is the experiment's most interesting result and we did not anticipate
it. Removing the shared-interface window cost exactly one acceptance case out of sixty
--- a near-null result on the metric we had chosen. But the agents told us why in their
own reports: deprived of their dependencies' published signatures, they did not fail,
they \emph{adapted}. One re-classified tokens by value because it could not know the
token-kind vocabulary; another constructed AST nodes by matching values onto whatever
field names the class actually declared; a third reached for \code{getattr} fallbacks
throughout. Each independently, and each without being asked to.

That behaviour is invisible to an acceptance suite and plain in the source, so we
measured it. Counting, with the AST rather than by pattern match, the constructs an
author writes when guessing at a collaborator's shape --- three-argument
\code{getattr}, \code{hasattr}, and handlers for the exceptions a wrong guess raises
--- gives \NSwDefShared{} per thousand lines with the window and \NSwDefNoshared{}
without it, a factor of \NSwDefRatio, alongside \NSwLinesRatio\% more code
(\NSwLinesShared{} lines versus \NSwLinesNoshared). The single agent, which owned every
module and therefore never had to guess, wrote \NSwSoloDef{} per thousand lines.

So the price of missing interface information is paid in runtime introspection rather
than in failed tests: the modules still compose, by each one defensively
accommodating whatever its neighbours turned out to be. That is a real cost and a
familiar one --- it is how human codebases decay when interfaces are informal --- and it
is exactly the cost a pass-rate metric cannot see. We offer
\emph{defensive constructs per thousand lines} as a cheap, deterministic, and
oracle-independent measure for multi-agent coding work generally; it needs no judge and
no test suite, only the source.

\begin{table}[t]
\centering\scriptsize
\begin{tabular}{lrrrr}
\toprule
module & deps & shared & no shared & ratio \\
\midrule
\multicolumn{5}{l}{\emph{precise specification}} \\
\quad\texttt{api.py} & 4 & 45 & 45 & 1.00 \\
\quad\texttt{engine.py} & 3 & 431 & 771 & \textbf{1.79} \\
\quad\texttt{errors.py} & 0 & 17 & 17 & 1.00 \\
\quad\texttt{functions.py} & 1 & 264 & 264 & 1.00 \\
\quad\texttt{nodes.py} & 0 & 179 & 179 & 1.00 \\
\quad\texttt{parser.py} & 3 & 374 & 711 & \textbf{1.90} \\
\quad\texttt{planner.py} & 2 & 520 & 537 & 1.03 \\
\quad\texttt{tokens.py} & 1 & 172 & 172 & 1.00 \\
\quad\emph{total} & & 2,002 & 2,696 & 1.35 \\
\addlinespace
\bottomrule
\end{tabular}

\caption{Code size per module with and without published interfaces. The cost falls
exactly on the modules that \emph{consume} interfaces: every module with no upstream
dependency is byte-identical, while the two modules with three dependencies each nearly
double.}
\label{tab:growth}
\end{table}

\paragraph{The cost lands exactly where the dependency graph says it should.} A global
``more code'' figure is weak evidence, because agents writing more code has many
explanations. \Cref{tab:growth} localises it, and the localisation is what makes it an
argument. Every module with no upstream dependency comes out byte-identical between the
two configurations --- \code{errors}, \code{nodes}, \code{tokens}, \code{functions} ---
which is the control: those authors had nothing to guess about, and they wrote the same
amount of code. The two modules with three dependencies each, \code{parser} and
\code{engine}, grew by 90\% and 79\%.

The agents said why. A reviewer that owns \code{engine} reported that, with nothing
published, it could not rely on the field names of \code{planner}'s \code{Plan} or of the
AST classes, so it dispatched on the class names the specification does fix, read fields
through a tolerant accessor, and \emph{did its own column binding, \code{*} expansion and
output naming as a fallback} --- work that belongs to \code{planner}, which the
specification explicitly forbids duplicating. This is a step past defensive reading: it
is responsibility duplication, the failure mode in which a module reimplements its
neighbour because it cannot trust it, and it is the mechanism by which informally
specified systems accrete.

It also sharpens what the two channels do. Both configurations had the precise
specification, so signatures were never in doubt; what the window carried was the
\emph{vocabulary and invariants behind} the signatures --- the token-\code{kind}
strings, the AST field names, what may be \code{NULL}. That is the information a caller
actually needs and a signature does not convey, and its absence cost 35\% more code
concentrated entirely among the consumers.

\paragraph{The pass-rate ablation was also confounded, and the confound is
instructive.} Beyond the adaptation above, there is a second reason the pass rate
barely moved. The
cause is a flaw in our experimental design: the broadcast specification already lists
each module's exported signatures, so \emph{the specification is itself a shared
interface}, and the window had nothing left to contribute. A rank could read
\code{tokenize(sql) -> list[Token]} out of the spec whether or not its author had
published anything.

The lesson is worth more than the ablation would have been:
\textbf{interface publication pays only to the extent that the specification
underdetermines the boundaries.} A precisely specified system does not need a
negotiation phase, and paying for one is waste; an underspecified system --- which is
the normal case, and the case in which eight agents each invent their neighbours'
signatures --- does. This also gives the mechanism a decision rule rather than a
recommendation: measure how much of your interface surface the specification pins, and
publish only the rest.

\paragraph{With an underspecified specification, interface publication pays
decisively.} We therefore ran a \emph{vague-specification} variant in which the module
table names responsibilities and import permissions but no signatures at all, leaving
every module boundary to be negotiated through the window --- the comparison we should
have designed first. It separates the two configurations cleanly on every axis we
measure.

With the window, the population reached \NVagueSharedRoundOne{} \emph{in a single
integration round} and stopped, in \NVagueSharedWall\,s using \NVagueSharedCalls{}
invocations and \$\NVagueSharedUsd. Without it, round one produced
\NVagueNosRoundOne{} and a second round was required to reach a full pass, taking
\NVagueNosWall\,s, \NVagueNosCalls{} invocations and \$\NVagueNosUsd. So publishing
interfaces saved \NVagueWallRatio$\times$ the wall time and \NVaguePriceRatio$\times$
the price, and --- the mechanism, not just the outcome --- it eliminated an entire
integrate-repair round. Defensive coupling differed by \NVagueDefRatio$\times$
(\NVagueSharedDef{} versus \NVagueNosDef{} per thousand lines) and code volume by
\NVagueLinesRatio$\times$ (\NVagueSharedLines{} versus \NVagueNosLines{} lines).

The per-module picture in \cref{tab:growth} changes shape in exactly the way the
mechanism predicts. Under the precise specification, only the consumers paid and the
dependency-free modules were byte-identical. Under the vague specification \emph{every}
module grows, including \code{errors} and \code{nodes}, which have no dependencies at
all --- because now every module must also be \emph{guessed at} by somebody, and an
author who does not know what its consumers will assume writes defensively in both
directions. Withholding the window costs 35\% more code when the specification is
precise and 62\% more when it is not.

This is the result the mechanism was designed to produce, and it took a redesigned
experiment to see it. The lesson for the protocol is the decision rule stated above;
the lesson for the methodology is that an ablation which removes a mechanism whose job
the specification is already doing measures nothing, and that we could not tell from
the pass rate alone that this was what had happened.

\paragraph{A directed collective on the wrong graph made peer review a no-op.} An
agent rank reported that the critiques it received named only a file it did not own.
It was right, and the cause is a mistake worth naming because the protocol makes it
easy to make. \code{review\_edges} yields \emph{(author, reviewer)} pairs, so a rank's
graph \emph{sources} are the authors whose code it reviews. Returning the critique with
a neighbourhood collective over that same graph sends it to the rank's
\emph{destinations} --- the ranks it sends its own code to --- which is a disjoint set.
Every author therefore received critiques of other people's modules, and the review
phase produced text that reached nobody who could act on it.

The fix is to return reviews over the \emph{transposed} graph, which is two lines. The
general lesson is that a neighbourhood collective on a directed topology has a
direction, and a request/response pair needs \emph{both} the graph and its transpose ---
a distinction MPI's \code{MPI\_Dist\_graph\_create} makes available (it takes sources
and destinations separately) and which we simply failed to use. Notably, nothing in the
run looked wrong: message counts, degrees and round counts were all correct, because
the traffic was well-formed and merely addressed to the wrong ranks. It was found by an
agent noticing the semantics, not by any test; there is now a test.

This compounds the confound above. In the precise-specification runs, the review phase
was inert \emph{and} the interface window was redundant, so two of the three
coordination mechanisms under study were doing nothing --- which is a further reason to
read that ablation as uninformative rather than negative. The vague-specification pair
runs with corrected routing.

\paragraph{Two further defects in our own measurement, both worth reporting.} The first was
in the oracle. One acceptance case tested case-insensitive \emph{keywords} using a
case-varied \emph{identifier}, which the specification says is case-sensitive; the
population implemented the specification correctly and the suite marked it wrong. An
oracle is a program and can be the defective party, which matters directly for
\cref{sec:ft}'s argument: a fault-tolerance scheme built on verification inherits the
reliability of its verifier, and we should not have claimed otherwise without saying
so.

The second was worse and less visible. The harness parsed the suite's JSON report by
slicing from the \emph{last} brace in its output --- which lands inside a nested
per-case object and never parses --- so every run was recorded as failing to import
with zero passes. Because that report is precisely what the harness scatters back to
the population as the definition of done, \textbf{the agents were told a build that
passed 58 of 59 cases had failed to import}, and spent their repair round on a
phantom. A bug in the plumbing \emph{around} an oracle is as damaging as a bug in the
oracle itself, and neither is visible from the population's behaviour.

We report the consequences rather than hiding them. Every configuration received the
same false signal, so the ablation comparison is between equally handicapped runs and
remains valid; the pass rates in \cref{tab:software} come from re-scoring the agents'
code offline against a single corrected oracle
(\code{scripts/reeval\_software.py}), which is free because the code is on disk. What
we cannot claim from this campaign is anything about \emph{repair} dynamics: the
second round was driven by a report that was not true, so the interesting question of
whether a population converges under honest failure feedback is untested here.
Notably, the population reached a full pass rate anyway --- it was robust to being
told, falsely, that its correct work had failed.

\subsection{Reduction fidelity}

\begin{table*}[t]
\centering\scriptsize
\emph{(not measured in this run)}

\caption{Fact retention of a semantic reduction. Each rank contributes a report of
uniquely identified factual items; retention is counted mechanically at the root.
``rank corr.'' is the correlation between a rank's index and the fraction of its
items that survived --- a positional unfairness \code{MPI\_SUM} cannot exhibit.}
\label{tab:fidelity}
\end{table*}

\Cref{tab:fidelity} tests the central prediction of \cref{sec:coll}: since all
algorithms perform the same number of operator applications, retention differences
must be attributable to fold depth.

\paragraph{The depth penalty is conditional, and the condition is capacity.} Our
first configuration --- eight ranks contributing four identified items each, merged
under a 900-token budget --- produced \emph{perfect} retention for all three binary
algorithms, with zero fabricated identifiers, whether the fold depth was 3 or 7. The
reason is straightforward once seen: thirty-two short items fit the output budget, so
the operator was never forced to discard anything, and an operator that discards
nothing is lossless at any depth.

This is a more useful claim than the one we set out to test. \textbf{Fold depth
penalises fidelity only when the reduction is capacity-bound}, so the first question
a harness author should ask is not which tree to use but whether their reduction is
saturated. If it is not, the shallowest tree is free and the fidelity argument is
moot; if it is, depth is the quantity to minimise. \Cref{tab:fidelity} therefore
reports a deliberately saturating configuration alongside the unsaturated one.

\paragraph{At equal quality, the tree wins outright on latency.} The unsaturated
configuration is still informative, because it isolates the cost axis with quality
held fixed at $1.0$: the binomial tree completed in \NFidTree\,s where the serial
chain took \NFidChain\,s and the flat reduction \NFidFlat\,s, all performing exactly
seven operator applications --- a $\NFidSpeedup\times$ latency reduction for free.
That flat is the \emph{slowest} despite needing one communication round is the cost
model's prediction confirmed: all seven applications execute serially at the root, so
its critical path is $7\alpha$ rather than $\lceil\log_2 8\rceil\alpha$. The
manager-summarises-everything pattern is the worst of the three on wall time and has
no compensating advantage on quality.

\subsection{Simulated scaling}

\begin{table}[t]
\centering\scriptsize
\begin{tabular}{llrrrrr}
\toprule
collective & algorithm & $p{=}8$ & $p{=}32$ & $p{=}128$ & $p{=}512$ & $p{=}1024$ \\
\midrule
allgather & \texttt{bruck} & 0 & 0 & 0 & 0 & 0 \\
allgather & \texttt{ring} & 0 & 0 & 0 & 2 & 3 \\
allreduce & \texttt{recursive\_doubling} & 159 & 290 & 451 & 596 & 682 \\
allreduce & \texttt{reduce\_bcast} & 159 & 290 & 451 & 596 & 682 \\
alltoall & \texttt{bruck} & 0 & 0 & 0 & 0 & 0 \\
alltoall & \texttt{pairwise} & 0 & 0 & 0 & 2 & 3 \\
barrier & \texttt{dissemination} & 0 & 0 & 0 & 0 & 0 \\
barrier & \texttt{linear} & 0 & 0 & 0 & 2 & 3 \\
bcast & \texttt{binomial} & 0 & 0 & 0 & 0 & 0 \\
bcast & \texttt{chain} & 0 & 0 & 0 & 2 & 3 \\
bcast & \texttt{flat} & 0 & 0 & 0 & 2 & 3 \\
reduce & \texttt{binomial} & 159 & 290 & 451 & 596 & 682 \\
reduce & \texttt{chain} & 329 & 1,416 & 6,104 & 24,246 & 48,737 \\
reduce & \texttt{flat} & 329 & 1,416 & 6,104 & 24,244 & 48,734 \\
\bottomrule
\end{tabular}

\caption{Simulated collective makespan (seconds) with parameters calibrated from the
measured runs, median of 32 trials. Latency is drawn log-normally, because agent
latency is right-skewed and a collective's completion time is a maximum over ranks.}
\label{tab:scalingsim}
\end{table}

No agent-systems paper can afford a 1024-rank experiment, and none should pretend
otherwise. HPC's honest answer is to build a discrete-event simulator of the
communication pattern, calibrate it against a real run at an affordable size, and
extrapolate~\cite{Hoefler10loggopsim}. \Cref{tab:scalingsim} is that
extrapolation, and \cref{tab:crossovers} the resulting algorithm-selection
crossovers --- the classical form of a collective-tuning result, now with a fidelity
column that has no MPI counterpart.

\begin{table}[t]
\centering\scriptsize
\begin{tabular}{lrrlll}
\toprule
collective & tokens & from $p$ & best (time) & best (volume) & best (fidelity) \\
\midrule
bcast & 64 & 2 & \texttt{binomial} & \texttt{binomial} & \texttt{binomial} \\
bcast & 512 & 2 & \texttt{binomial} & \texttt{binomial} & \texttt{binomial} \\
bcast & 4,096 & 2 & \texttt{binomial} & \texttt{binomial} & \texttt{binomial} \\
bcast & 32,768 & 2 & \texttt{binomial} & \texttt{binomial} & \texttt{binomial} \\
reduce & 64 & 2 & \texttt{binomial} & \texttt{binomial} & \texttt{binomial} \\
reduce & 512 & 2 & \texttt{binomial} & \texttt{binomial} & \texttt{binomial} \\
reduce & 4,096 & 2 & \texttt{binomial} & \texttt{binomial} & \texttt{binomial} \\
reduce & 32,768 & 2 & \texttt{binomial} & \texttt{binomial} & \texttt{binomial} \\
reduce & 64 & 4 & \texttt{kary} & \texttt{binomial} & \texttt{kary} \\
reduce & 512 & 4 & \texttt{kary} & \texttt{binomial} & \texttt{kary} \\
reduce & 4,096 & 4 & \texttt{kary} & \texttt{binomial} & \texttt{kary} \\
reduce & 32,768 & 4 & \texttt{kary} & \texttt{binomial} & \texttt{kary} \\
allreduce & 64 & 2 & \texttt{recursive\_doubling} & \texttt{recursive\_doubling} & \texttt{recursive\_doubling} \\
allreduce & 512 & 2 & \texttt{recursive\_doubling} & \texttt{recursive\_doubling} & \texttt{recursive\_doubling} \\
allreduce & 4,096 & 2 & \texttt{recursive\_doubling} & \texttt{recursive\_doubling} & \texttt{recursive\_doubling} \\
allreduce & 32,768 & 2 & \texttt{recursive\_doubling} & \texttt{recursive\_doubling} & \texttt{recursive\_doubling} \\
allgather & 64 & 2 & \texttt{bruck} & \texttt{bruck} & \texttt{bruck} \\
allgather & 512 & 2 & \texttt{bruck} & \texttt{bruck} & \texttt{bruck} \\
allgather & 4,096 & 2 & \texttt{bruck} & \texttt{bruck} & \texttt{bruck} \\
allgather & 32,768 & 2 & \texttt{bruck} & \texttt{bruck} & \texttt{bruck} \\
alltoall & 64 & 2 & \texttt{bruck} & \texttt{bruck} & \texttt{bruck} \\
alltoall & 512 & 2 & \texttt{bruck} & \texttt{bruck} & \texttt{bruck} \\
alltoall & 4,096 & 2 & \texttt{bruck} & \texttt{bruck} & \texttt{bruck} \\
alltoall & 32,768 & 2 & \texttt{bruck} & \texttt{bruck} & \texttt{bruck} \\
barrier & 64 & 2 & \texttt{dissemination} & \texttt{central} & \texttt{central} \\
barrier & 512 & 2 & \texttt{dissemination} & \texttt{central} & \texttt{central} \\
barrier & 4,096 & 2 & \texttt{dissemination} & \texttt{central} & \texttt{central} \\
barrier & 32,768 & 2 & \texttt{dissemination} & \texttt{central} & \texttt{central} \\
barrier & 64 & 4 & \texttt{central} & \texttt{central} & \texttt{central} \\
barrier & 512 & 4 & \texttt{central} & \texttt{central} & \texttt{central} \\
barrier & 4,096 & 4 & \texttt{central} & \texttt{central} & \texttt{central} \\
barrier & 32,768 & 4 & \texttt{central} & \texttt{central} & \texttt{central} \\
\bottomrule
\end{tabular}

\caption{Algorithm recommendations under three objectives. Where the time-optimal
and fidelity-optimal choices differ, the harness author has a decision to make that
MPI never presents.}
\label{tab:crossovers}
\end{table}

\section{Discussion and limitations}
\label{sec:limitations}

\paragraph{Scale.} Our real-agent populations are single-digit, because the
concurrency limit of our agent host bound them, and because the campaign structure
--- one persistent worker pool serving a whole ablation ladder --- trades population
size for the ability to run many configurations against the same agents. Larger
populations are reported only as calibrated simulation and labelled as such. The
protocol-level results are size-independent and validated to $p=32$.

\paragraph{The insensitive metric.} As noted in Q3, our entity-consistency metric
saturates on famous proper nouns. A more sensitive replication should ask each rank
to report renderings for \emph{every} proper noun it used rather than a
pre-registered list, and score consistency over terms reported by two or more
ranks; rare names diverge where famous ones do not. We report the saturated result
rather than substituting the metric after seeing the data.

\paragraph{An exact operator can impose a locally wrong global decision, and our
metric rewards it.} A worker rank surfaced a case we had not anticipated. The
\code{UNION} glossary bound the source term \emph{pounds} to the rendering for pounds
sterling, correctly for most of the book; in one section the word denotes body weight,
and the rank --- correctly following a glossary the harness declared binding ---
produced a semantically wrong sentence. The section that received no glossary rendered
it naturally.

This is a real limitation of the design, not an implementation slip. An exact,
order-independent merge over a term map presupposes that a source term has one
correct target rendering, which polysemy falsifies; \code{UNION}'s conflict retention
handles two ranks \emph{disagreeing} about a term but not one term having two correct
senses. The fix is a contract change --- key the glossary on (term, sense) and let a
rank declare which sense it saw --- and it makes the shared state larger and the
agreement weaker, which is exactly the trade the design should have surfaced.

Worse for our evaluation: \textbf{the consistency metric scores this failure as a
success}, because it counts identical renderings across sections and a globally
imposed wrong rendering is maximally consistent. A metric that rewards agreement
cannot detect agreement on something false. That is a specific instance of the general
hazard in \cref{sec:ft} --- verification inherits the reliability of its verifier ---
and it means our perfect consistency scores should be read as evidence about
\emph{coordination}, which is what they measure, and not about translation quality,
which they do not.

\paragraph{We committed the lost-update bug in the harness that runs the experiments.}
Two campaigns sharing a campaign directory ran concurrently; the second activated its
job, then the first finished and cleared the shared pointer with a blind write,
stranding the second's worker pool with nothing to poll. A reduction sat waiting on a
rank that could no longer find its work. That is exactly the two-writers-no-
synchronisation failure that \cref{sec:rma} is about, and the remedy is the one that
section prescribes --- clear the cell only if it still holds your value, a
compare-and-swap. We record it because it is the paper's own argument arriving as a
self-inflicted wound: the discipline is easy to state, easy to skip, and its violation
fails silently rather than loudly. A protocol that makes the safe operation available
does not make anyone use it, which is an argument for defaults, not for documentation.

\paragraph{Shared mutable code is shared state the protocol does not cover.} One
worker crashed with a type error inside the runtime, and the cause was that we edited
the package --- installed in editable mode --- while a live population was executing
against it, so that worker observed a half-completed refactor. The protocol keeps
protocol \emph{state} outside the agents and durable, which is the design's central
move, and it does nothing at all about the \emph{code} being shared and mutable. A
production deployment should pin an immutable runtime version per job and have workers
report a mismatch; ours did not, and the honest description of the failure is that we
hot-patched a running job.

\paragraph{One model family.} All agent ranks were served by one model family, so we
cannot separate protocol effects from model idiosyncrasies, and the correlated-failure
caution in \cref{sec:ft} applies to our own replication results.

\paragraph{Token counting.} Where provider usage was unavailable we used a
documented estimator; every conclusion depending on exact counts is drawn from
ratios rather than absolute values.

\paragraph{Fabric as a fate-sharing point.} The reference fabric is a single durable
store, so \code{agree} routes agreement through a reliable service rather than
circumventing FLP. The specification states that requirement rather than hiding it;
a distributed fabric would need real consensus.

\paragraph{What we would change.} The context budget is currently a scalar. Real
agents have structured contexts --- system prompt, tool definitions, conversation,
retrieved documents --- with very different eviction economics, and a protocol that
modelled the structure could make better admission decisions.

\section{Related work}
\label{sec:related}

\paragraph{MPI and its lineage.} We build directly on the MPI standard and its
algorithmic
literature~\cite{mpiforum94,mpi31,mpi40,ThakurRabenseifnerGropp05,Chan07,Bruck97,Rabenseifner04},
its one-sided
design~\cite{hoefler15-rma,dinan16-rma,gerstenberger14}, and its fault-tolerance
line~\cite{fagg00-ftmpi,bouteiller22}. Our departures are catalogued in
\cref{sec:intro}.

\paragraph{Other parallel models.} BSP's superstep discipline~\cite{valiant90bsp} is
the ancestor of our barrier-separated phases; Pregel applies it to
graphs~\cite{malewicz10pregel}. The actor
model~\cite{hewitt73actor,agha86actors} and Erlang/OTP~\cite{armstrong03erlang}
supply supervision, which MPI lacks and which we adopt wholesale; Orleans supplies
the virtual-actor idea behind rank incarnations~\cite{bernstein14orleans}. CSP
contributes the rendezvous channel~\cite{hoare78csp}.
MapReduce and Spark show how restricting expressiveness buys
recoverability~\cite{dean04mapreduce,zaharia12rdd}; we make the opposite trade and
pay for it with an explicit failure model. Task-graph
runtimes~\cite{bauer12legion,moritz18ray} and pipeline
parallelism~\cite{huang19gpipe} inform our topology work.

\paragraph{Agent frameworks.} AutoGen, MetaGPT, ChatDev, CAMEL, LangGraph and CrewAI
each supply a coordination model, and each is a framework rather than a
protocol: they own the control flow, and a harness written against one cannot be
composed with
another~\cite{wu24autogen,hong24metagpt,qian24chatdev,li23camel,langgraph,crewai}.
None provides scoped communication contexts, collectives with selectable algorithms,
one-sided state with defined concurrency semantics, or a specified behaviour when a
participant disappears. Durable-execution systems~\cite{temporal} give reliable
step execution but no communication algebra.

\paragraph{Agent protocols.} MCP standardises agent-to-tool
access~\cite{mcpspec2026}; A2A standardises task delegation across
organisations~\cite{a2adefs}; ANP, ACP and Coral address discovery and
routing~\cite{anpwp,coral}. Surveys confirm the
division~\cite{yangsurvey,ehteshamsurvey}. These are complementary and orthogonal:
they connect \emph{systems}, whereas \ampi{} is what one writes a system
\emph{with}. Two details make the orthogonality sharper than it first appears. MCP's
current revision moved deliberately \emph{away} from shared state, making requests
stateless and self-contained with per-request capability negotiation, specifically to
avoid stateful load balancing~\cite{mcpspec2026,mcpsep2322}; it is a tool protocol
hardening into statelessness, not drifting toward agent-to-agent. And A2A's
\code{contextId} is a correlation identifier rather than a membership set: its unit
of coordination is the bilateral task, with retry and circuit-breaking explicitly
left to the client~\cite{a2adefs,a2alifeoftask}. Neither has a group, and without a
group there is no collective.

Classical agent communication languages --- KQML, FIPA-ACL, the Contract Net
Protocol, blackboard architectures, Linda tuple
spaces~\cite{kqml,fipaacl,contractnet,hearsay,carriero89-linda,jade} --- are the true
ancestors, and modern protocols are rediscovering them; notably, Linda's tuple space
is a shared-memory abstraction with no epochs, which is precisely the gap
\cref{sec:rma} fills.

\paragraph{The closest prior work, and the objection it raises.} Mieczkowski et
al.~\cite{lmteams26} is simultaneously our strongest motivation and the paper a
reviewer will ask about. It validates the distributed-systems analogy empirically ---
Amdahl bounds, sub-unity speedup for decentralised teams, concurrent-write conflicts
as a dominant failure --- and it is the work we would cite to justify the whole
enterprise. It produces analysis and design guidance, not an interface. Our claim is
therefore narrower and more specific than ``agent teams are distributed systems'':
it is that a particular thirty-year-old interface transfers, that three named
modifications are forced, and that the modifications are measurable. We are also
aware of at least one independent design sketch proposing MPI-shaped primitives for
agents (communicators, deterministic reduce, one-sided put/get, shrink and agree) in
a project issue tracker rather than a publication; we claim priority on none of the
vocabulary, only on the specification, the implementation, and the measurements.

\paragraph{Reduction cost in MPI's own literature.} Two results deserve particular
credit because they anticipate our concerns from inside MPI. The commutativity flag
on \code{MPI\_Op\_create} is already an algorithm-selection gate rather than
documentation, which is exactly the mechanism we generalise to a three-valued
associativity declaration. And Tr\"aff's recent scan work is explicitly motivated by
operators whose \emph{application} is expensive, minimising the number of operator
applications subject to a bound on rounds~\cite{traff25exscan} --- the closest thing
in the MPI literature to an algorithm designed for an LLM cost model, and the same
objective our \code{kary} reduction optimises. On reproducibility, the contrast
between Intel's conditional numerical reproducibility (reproducible only at a fixed
thread count) and ReproBLAS's binned accumulator (reproducible regardless of order or
process count)~\cite{IntelCNR,ReproBLAS,Demmel15} makes the design argument for us:
fix the \emph{operator}, not the order --- which is why \cref{sec:coll} pushes
harnesses toward exact, order-independent operators like \code{UNION} wherever the
combination is genuinely set-like.

\paragraph{Agent failure and context.} The MAST taxonomy attributes multi-agent
failures largely to specification and inter-agent
coordination~\cite{mast,mastneurips}; practitioner reports converge on context
fragmentation~\cite{cognitiondont,anthropicmas,anthropicctxeng,lcbenchmark}.
Long-context degradation is
measured~\cite{liu24lostinmiddle,labanlost,chromarot}, as are error cascades and
diversity collapse~\cite{diversitycollapse}. Agent-OS work treats context as a
managed resource~\cite{memgpt,aios}; serving systems schedule agent
workloads~\cite{parrot,autellix,teola}. Our contribution relative to all of these is
to make context an \emph{accounted protocol resource} with a transport-mode policy,
rather than a memory-management strategy inside one agent.

\paragraph{Multi-agent translation and coding.} Multi-agent literary translation has
been studied~\cite{transagentstacl,transagentsdemo}, and coding benchmarks are
standard~\cite{jimenez24swebench}. We use both as \emph{vehicles} for measuring
protocol mechanisms, not as targets; our harnesses are deliberately simple so that
ablations isolate the coordination primitive rather than prompt engineering.

\section{Conclusion}

The multi-agent field is rediscovering, application by application, the problems the
MPI Forum solved once for parallel computing. We have shown that MPI's core
abstractions --- communicators, typed transfers with projections, collectives with
selectable algorithms, one-sided state with epochs, and failure mitigation ---
transfer to agent populations, and that three modifications are forced and
consequential: transport mode must be semantically visible because context is the
scarce buffer; reduction operators must declare their algebra because fidelity
decays in fold depth; and fault tolerance must be built around verification because
the dominant failure is silent corruption rather than crash. What makes the
transplant work is one decision that MPI never had to make and that agent frameworks
never make: artifacts are immutable and content-addressed, which is why broadcast can
be logarithmic and lossless while reduction cannot.

We do not claim \ampi{}'s spelling is right. MPI's was not right either in 1994, and
the value of the Forum was the process, not the first draft. We claim that the
\emph{shape} of the problem is the same one message passing already solved, and that
a field currently shipping one coordination mechanism per project should look at what
happened after 1994.

\balance
\bibliographystyle{plain}
\bibliography{agentmpi}

\end{document}
